% Explanation certificate for heart-s2-c3 (genome 1a5255ec-4f4d-4ff4-916a-1b5c4622ccc9).
% Generated by explaneat export (schema 1.4.0) from the JSON record alone;
% content sha256 6b36c4ae4f91161350a650f7da79e66ec93b2fa97254ec7aeab35bbb6828150b.
% Drafter recorded in the JSON: ai_agent (claude). Do not edit by hand: re-render with
%   uv run python -m explaneat export render --from heart-s2-c3.json
% TikZ modules are input from \explaneatcertdir{<label>} (empty by default: same
% directory; index.tex sets it to the genome's subdirectory).
\providecommand{\explaneatcertdir}[1]{}
\begin{drafted}[appendix: heart-s2-c3 explanation certificate]
\section{Explanation certificate: heart-s2-c3}\label{app:heart-s2-c3:sec}
This section is the explanation certificate of the heart-s2-c3 network. It is rendered from the canonical JSON record \texttt{heart-s2-c3.json} (schema 1.4.0, SHA-256 \texttt{6b36c4ae4f911613}\ldots), which also holds every evidence payload, the full graphs with weights, and the explanation record. \Cref{app:heart-s2-c3:summary} summarises the network and its coverage.

\begin{table}[htbp]
\centering\small
\begin{tabularx}{\linewidth}{@{}lX@{}}
\toprule
Dataset & Heart Disease (binary, prepared) \\
Experiment & 10seed-heart-s2, seed 2 \\
Genome & \texttt{1a5255ec-4f4d-4ff4-916a-1b5c4622ccc9} (NEAT key 9600, generation 39) \\
Explanation & \texttt{6565cffd-9c19-4ae2-b635-a2e35f3d485a} (34 operations) \\
Train AUC / accuracy & 0.9381 / 0.8884 (n = 242) \\
Test AUC / accuracy & 0.8987 / 0.8197 (n = 61) \\
Output activation & sigmoid \\
Phenotype & 30 nodes (26 inputs, 3 hidden, 1 output), 26 enabled edges \\
Explained model state & 36 nodes, 26 enabled edges \\
Refactoring operations & split\_\allowbreak{}node 4 \\
Function preservation & max $|\Delta\,\mathrm{output}|$ = $2.22\times 10^{-16}$ on the training split \\
Annotations & 4 (3 leaf, 1 composition) \\
Structural coverage & 1 (35 of 35 candidate nodes; 9 Null) \\
Compositional coverage & 1 \\
Uncovered nodes / edges & 0 / 0 \\
Missing compositions & 0 \\
Full coverage & yes \\
Drafted by & AI-assisted draft \\
Verified by & unverified \\
Record hash (SHA-256) & \texttt{6b36c4ae4f911613}\ldots \\
\bottomrule
\end{tabularx}
\caption{Explanation certificate summary: heart-s2-c3.}
\label{app:heart-s2-c3:summary}
\end{table}

\paragraph{Drafter's summary.} Mostly logistic regression: 16 direct edges (logit SD 3.67; train AUC 0.934 alone) plus three additive ReLU gates (logit SD 0.42) to 0.938. All three gates are switched by non-anginal chest pain (cp:3), which has no direct edge: a protective gate ReLU(-0.35 + 1.08 cp:3 + 0.68 thalach) read out -0.51, open for 97\% of non-anginal-pain patients, and a smaller low-cholesterol gate for the same patients (-0.24), plus a near-null oldpeak hinge for everyone else. Age is not used at all. Unlike the 10-seed heart-s2 record (key 9083, an age/\allowbreak{}heart-rate-reserve gate and oldpeak hinges), this genome's only nonlinear effect is a chest-pain-type-specific level and thalach/cholesterol slope.

\paragraph{What resisted explanation.} Nothing resisted coverage. The ST-depression gate for patients without non-anginal pain is a near-null part (readout -0.018, SD 0.009 logit units; ablating it raises train AUC by 0.0001) and is annotated as such. Counter-directional signs reported, not explained: cp:2 +0.27 and fbs:1 -0.21 in the linear core, extra protection for non-anginal-pain patients with lower cholesterol, and the gate that slightly lowers risk with more ST depression.

\subsection{Plain-language summary}\label{app:heart-s2-c3:plain}
\noindent{\small\textit{AI-assisted draft for review (unverified). Each sentence cites the claims (C01, C02, \ldots) that back it; every cited claim was computed on the model and the training split and passed its automatic checks. Magnitudes are changes in predicted risk for a typical patient; confidence words follow the test-split check and n. The text describes what the model does, not the world.}}\par\smallskip
\paragraph{Summary (drafted).} For a typical patient (a 56-year-old man with no chest pain symptoms and median test results) the model predicts a risk of heart disease of about 30\%; it adds up most of its inputs independently, and three parts that all read chest pain type add to that sum: one acts on 76 of 242 training patients through maximum heart rate, one on 61 through cholesterol, and one barely changes the predictions.\,{\scriptsize\textcolor{NavyBlue}{[C29, C34, C40, C50]}} These add up independently, and these lower the risk: through this sum alone, compared with no chest pain symptoms, typical angina takes a typical patient's predicted risk from about 30\% to about 2\%, non-anginal chest pain to about 7\% and atypical angina to about 13\% (chest pain type also acts through the parts below); being a woman lowers it to about 4\% and fasting blood sugar above 120 mg/dL to about 19\%; and through this sum alone a higher maximum heart rate lowers it steadily, from about 36\% at 118 to about 27\% at 175 beats per minute (it also acts through the parts below).\,{\scriptsize\textcolor{NavyBlue}{[C05, C14, C16, C10, C18, C19]}} Each extra vessel coloured by fluoroscopy adds about the same, from about 30\% at 0 to about 96\% at 2 vessels, flattening out near the top because the predicted risk cannot go much higher; a flat rather than upsloping ST segment raises it to about 80\%, higher resting blood pressure raises it steadily from about 20\% at 110 mm Hg to about 46\% at 155 mm Hg, and a reversible defect on the thallium scan raises it to about 55\%.\,{\scriptsize\textcolor{NavyBlue}{[C01, C02, C11, C12]}} Through this sum alone, more ST depression on exercise raises it steadily from about 23\% at 0 mm to about 53\% at 2.8 mm, and higher cholesterol raises it steadily from about 21\% at 193 mg/dL to about 48\% at 308 mg/dL; both also act through the parts below.\,{\scriptsize\textcolor{NavyBlue}{[C03, C07]}} Some rarer results also count (the model has little data on this, n = 15 to 17): a downsloping ST segment raises the risk to about 54\% and a fixed defect on the thallium scan lowers it to about 12\%, while signs of left ventricular hypertrophy on the resting ECG raise it from about 30\% to about 33\%.\,{\scriptsize\textcolor{NavyBlue}{[C13, C15, C21]}} This part reads chest pain type and maximum heart rate, and acts on 76 of 242 training patients, mostly patients with non-anginal chest pain, 15 of whom had heart disease; for the typical patient it acts on, the predicted risk is about 5\% without it and about 2\% with it.\,{\scriptsize\textcolor{NavyBlue}{[C34]}} Through it, non-anginal chest pain rather than any other chest pain type takes a typical patient's predicted risk from about 30\% to about 17\%.\,{\scriptsize\textcolor{NavyBlue}{[C30]}} Across the whole model, a higher maximum heart rate lowers the risk slowly at first, then faster, and more with non-anginal chest pain: from about 7\% at 118 to about 2\% at 175 beats per minute then, flattening out near the bottom because the predicted risk cannot go much lower, against about 36\% to about 27\% for other chest pain types.\,{\scriptsize\textcolor{NavyBlue}{[C32]}} This part reads chest pain type and cholesterol, and acts on 61 of 242 training patients, all of them with non-anginal chest pain, 13 of whom had heart disease; for the typical patient it acts on, the predicted risk is about 3\% without it and about 2\% with it.\,{\scriptsize\textcolor{NavyBlue}{[C40]}} Through it, for patients with non-anginal chest pain, higher cholesterol raises the predicted risk steadily up to about 282 mg/dL, where the part switches off, then no further: for the typical patient it acts on, from about 2\% at 193 mg/dL to about 3\% at 308 mg/dL.\,{\scriptsize\textcolor{NavyBlue}{[C37, C39]}} Across the whole model, higher cholesterol raises the risk in an uneven way from about 2\% to about 7\% with non-anginal chest pain, against about 21\% to about 48\% for other chest pain types.\,{\scriptsize\textcolor{NavyBlue}{[C38]}} This part reads chest pain type, maximum heart rate and ST depression on exercise, and barely changes the predictions (by at most 0.5 percentage points).\,{\scriptsize\textcolor{NavyBlue}{[C50]}} It switches on when ST depression on exercise minus about 0.019 $\times$ the maximum heart rate exceeds about -1.71, which never happens with non-anginal chest pain; 73 of 242 training patients are on, and 79\% of them had heart disease, against 29\% of the others.\,{\scriptsize\textcolor{NavyBlue}{[C52]}}
\paragraph{The whole model.} The model adds up its inputs into one score: the vessels coloured (each extra vessel adding about the same), the ST segment, sex and chest pain type each move the predicted risk on their own, and three parts add to that sum, two of which lower the risk further for patients with non-anginal chest pain.\,{\scriptsize\textcolor{NavyBlue}{[C29, C01, C02, C10, C09]}} \textit{Evidence.} It is made up of the inputs that feed the score directly and three parts, the non-anginal chest pain parts with heart rate and with cholesterol and an ST depression gate (each explained below); the score is their sum, which was checked on all 242 training patients.\,{\scriptsize\textcolor{NavyBlue}{[C29]}} Changing one input at a time across the whole model, non-anginal chest pain rather than no chest pain symptoms takes a typical patient's predicted risk from about 30\% to about 3\%, against about 7\% through the sum alone.\,{\scriptsize\textcolor{NavyBlue}{[C09, C14]}} Across the whole model, a higher maximum heart rate lowers it slowly at first, then faster, from about 36\% at 118 to about 27\% at 175 beats per minute, where through the sum alone it falls steadily over the same range.\,{\scriptsize\textcolor{NavyBlue}{[C20, C19]}} Across the whole model, more ST depression raises it steadily from about 23\% to about 53\%, and higher cholesterol steadily from about 21\% to about 48\%.\,{\scriptsize\textcolor{NavyBlue}{[C04, C08]}} It ignores age and chest pain brought on by exercise.\,{\scriptsize\textcolor{NavyBlue}{[C27, C28]}}
\paragraph{Non-anginal chest pain with heart rate.} This part lowers the predicted risk for patients with non-anginal chest pain, and for them a higher maximum heart rate lowers it further, slowly at first, then faster.\,{\scriptsize\textcolor{NavyBlue}{[C30, C32]}} \textit{Evidence.} It acts on 76 of 242 training patients, mostly patients with non-anginal chest pain, 15 of whom had heart disease; for the typical patient it acts on, the predicted risk is about 5\% without it and about 2\% with it.\,{\scriptsize\textcolor{NavyBlue}{[C34]}} Through it, non-anginal chest pain rather than any other chest pain type takes a typical patient's predicted risk from about 30\% to about 17\%.\,{\scriptsize\textcolor{NavyBlue}{[C30]}} Through it, a higher maximum heart rate lowers the risk in an uneven way, for the typical patient it acts on from about 4\% at 118 to about 2\% at 175 beats per minute, flattening out near the bottom because the predicted risk cannot go much lower.\,{\scriptsize\textcolor{NavyBlue}{[C31]}} Across the whole model, a higher maximum heart rate lowers the risk slowly at first, then faster, from about 7\% to about 2\% with non-anginal chest pain, against about 36\% to about 27\% for other chest pain types.\,{\scriptsize\textcolor{NavyBlue}{[C32]}} Its formula reproduces the part exactly on all 242 training patients.\,{\scriptsize\textcolor{NavyBlue}{[C35]}}
\paragraph{Non-anginal chest pain with cholesterol.} This part lowers the predicted risk a little further for patients with non-anginal chest pain whose cholesterol stays below about 282 mg/dL.\,{\scriptsize\textcolor{NavyBlue}{[C40, C39]}} \textit{Evidence.} It acts on 61 of 242 training patients, all of them with non-anginal chest pain, 13 of whom had heart disease; for the typical patient it acts on, the predicted risk is about 3\% without it and about 2\% with it.\,{\scriptsize\textcolor{NavyBlue}{[C40]}} Through it, higher cholesterol raises the predicted risk steadily up to about 282 mg/dL, where the part switches off, then no further: for the typical patient it acts on, from about 2\% at 193 mg/dL to about 3\% at 308 mg/dL.\,{\scriptsize\textcolor{NavyBlue}{[C37, C39]}} Through it, non-anginal chest pain rather than any other chest pain type takes a typical patient's predicted risk from about 30\% to about 29\%.\,{\scriptsize\textcolor{NavyBlue}{[C36]}} Across the whole model, higher cholesterol raises the risk in an uneven way from about 2\% to about 7\% with non-anginal chest pain, against about 21\% to about 48\% for other chest pain types.\,{\scriptsize\textcolor{NavyBlue}{[C38]}} Its formula reproduces the part exactly on all 242 training patients.\,{\scriptsize\textcolor{NavyBlue}{[C41]}}
\paragraph{ST depression gate with little effect.} This part barely changes the predictions.\,{\scriptsize\textcolor{NavyBlue}{[C50]}} \textit{Evidence.} It switches on when ST depression on exercise minus about 0.019 $\times$ the maximum heart rate exceeds about -1.71, which never happens with non-anginal chest pain; 73 of 242 training patients are on, and 79\% of them had heart disease, against 29\% of the others.\,{\scriptsize\textcolor{NavyBlue}{[C52]}} It changes the predicted risk by at most 0.5 percentage points: through it, ST depression on exercise and maximum heart rate barely move a typical patient's predicted risk, which stays at about 30\%.\,{\scriptsize\textcolor{NavyBlue}{[C50, C42, C43]}} Its formula reproduces the part exactly on all 242 training patients.\,{\scriptsize\textcolor{NavyBlue}{[C51]}}
\paragraph{Worth checking against what is known.} Worth checking: going from 0 to 2 vessels coloured moves a typical patient's predicted risk from about 30\% to about 96\%.\,{\scriptsize\textcolor{NavyBlue}{[C01]}} Worth checking: patients with typical angina get the lowest predicted risk of all chest pain types (about 2\%), below non-anginal chest pain (about 3\%), atypical angina (about 13\%) and no chest pain symptoms (about 30\%).\,{\scriptsize\textcolor{NavyBlue}{[C06, C09, C17]}} Worth checking: the model treats fasting blood sugar above 120 mg/dL as lowering the risk of heart disease (for a typical patient from about 30\% to about 19\%).\,{\scriptsize\textcolor{NavyBlue}{[C18]}} Worth checking: the model treats a fixed defect on the thallium scan as lowering the risk (for a typical patient from about 30\% to about 12\%), from only 15 training patients.\,{\scriptsize\textcolor{NavyBlue}{[C15]}} Worth checking: the model ignores age.\,{\scriptsize\textcolor{NavyBlue}{[C27]}} Worth checking: with non-anginal chest pain, one part lowers the predicted risk for cholesterol below about 282 mg/dL, so cholesterol moves the risk from about 2\% to about 7\% for these patients, against about 21\% to about 48\% for others.\,{\scriptsize\textcolor{NavyBlue}{[C38, C39]}}
\paragraph{Claims.}
\par{\tiny\noindent
\begin{tabularx}{\linewidth}{@{}l>{\raggedright\arraybackslash}X>{\raggedright\arraybackslash}p{0.11\linewidth}>{\raggedright\arraybackslash}p{0.16\linewidth}>{\raggedright\arraybackslash}p{0.05\linewidth}>{\raggedright\arraybackslash}p{0.15\linewidth}l>{\raggedright\arraybackslash}p{0.08\linewidth}@{}}
\toprule
Claim & What is checked & Risk, typical patient & Estimate (log-odds; OR); range across patients & n & Test check & Verified & Words \\
\midrule
C01 & direction: number of vessels coloured by fluoroscopy: 0 vessels to 2 vessels (whole model) \newline\textit{shape: steady (equal steps); 100\% of patients same; verified; risk saturates near 100\%} & 30\% to 96\% & log-odds 3.99, OR 54.1; range across patients 3.99 to 3.99 & 242 & train; test unclear (AUC gain 0.039 [-0.034, 0.120], n=61) & \textcolor{OliveGreen}{yes} & on the training data only \\
C02 & direction: slope of the ST segment at peak exercise: an upsloping ST segment to a flat ST segment (whole model) \newline\textit{shape: categorical; highest to lowest: a flat ST segment, a downsloping ST segment, an upsloping ST segment} & 30\% to 80\% & log-odds 2.20, OR 9.04; range across patients 2.20 to 2.20 & 115, 110 & train; test disagrees (AUC gain -0.005 [-0.050, 0.037], n=61) & \textcolor{OliveGreen}{yes} & on the training data only \\
C03 & direction: ST depression on exercise: 0.0 mm to 2.8 mm (direct contribution only, not through any part) \newline\textit{shape: steady; 100\% of patients same; verified} & 23\% to 53\% & log-odds 1.32, OR 3.74; range across patients 1.32 to 1.32 & 242 & train; test unclear (AUC gain 0.011 [-0.01, 0.038], n=61) & \textcolor{OliveGreen}{yes} & on the training data only \\
C04 & direction: ST depression on exercise: 0.0 mm to 2.8 mm (whole model) \newline\textit{shape: steady; 100\% of patients same; verified} & 23\% to 53\% & log-odds 1.31, OR 3.69; range across patients 1.30 to 1.32 & 242 & train; test unclear (AUC gain 0.011 [-0.01, 0.038], n=61) & \textcolor{OliveGreen}{yes} & on the training data only \\
C05 & direction: chest pain type: no chest pain (asymptomatic) to typical angina (direct contribution only, not through any part) \newline\textit{shape: categorical; highest to lowest: no chest pain (asymptomatic), atypical angina, non-anginal chest pain, typical angina} & 30\% to 2\% & log-odds -3.11, OR 0.04; range across patients -3.11 to -3.11 & 108, 22 & train; test agrees (AUC gain 0.038 [0.011, 0.072], n=61) & \textcolor{OliveGreen}{yes} & clearly \\
C06 & direction: chest pain type: no chest pain (asymptomatic) to typical angina (whole model) \newline\textit{shape: categorical; highest to lowest: no chest pain (asymptomatic), atypical angina, non-anginal chest pain, typical angina} & 30\% to 2\% & log-odds -3.11, OR 0.04; range across patients -3.11 to -3.11 & 108, 22 & train; test agrees (AUC gain 0.055 [0.014, 0.101], n=61) & \textcolor{OliveGreen}{yes} & clearly \\
C07 & direction: serum cholesterol: 193 mg/dL (5.0 mmol/L) to 308 mg/dL (8.0 mmol/L) (direct contribution only, not through any part) \newline\textit{shape: steady; 100\% of patients same; verified} & 21\% to 48\% & log-odds 1.26, OR 3.54; range across patients 1.26 to 1.26 & 242 & train; test disagrees (AUC gain -0.012 [-0.044, 0.013], n=61) & \textcolor{OliveGreen}{yes} & on the training data only \\
C08 & direction: serum cholesterol: 193 mg/dL (5.0 mmol/L) to 308 mg/dL (8.0 mmol/L) (whole model) \newline\textit{shape: steady; 100\% of patients same; verified} & 21\% to 48\% & log-odds 1.31, OR 3.71; range across patients 1.26 to 1.43 & 242 & train; test disagrees (AUC gain -0.012 [-0.044, 0.013], n=61) & \textcolor{OliveGreen}{yes} & on the training data only \\
C09 & direction: chest pain type: no chest pain (asymptomatic) to non-anginal chest pain (whole model) \newline\textit{shape: categorical; highest to lowest: no chest pain (asymptomatic), atypical angina, non-anginal chest pain, typical angina} & 30\% to 3\% & log-odds -2.51, OR 0.08; range across patients -2.96 to -2.01 & 108, 70 & train; test agrees (AUC gain 0.055 [0.014, 0.101], n=61) & \textcolor{OliveGreen}{yes} & clearly \\
C10 & direction: sex: male to female (whole model) \newline\textit{shape: categorical; highest to lowest: male, female} & 30\% to 4\% & log-odds -2.33, OR 0.10; range across patients -2.33 to -2.33 & 159, 83 & train; test unclear (AUC gain 0.025 [-0.011, 0.071], n=61) & \textcolor{OliveGreen}{yes} & on the training data only \\
C11 & direction: resting blood pressure: 110 mm Hg to 155 mm Hg (whole model) \newline\textit{shape: steady; 100\% of patients same; verified} & 20\% to 46\% & log-odds 1.22, OR 3.39; range across patients 1.22 to 1.22 & 242 & train; test unclear (AUC gain 0.004 [-0.019, 0.032], n=61) & \textcolor{OliveGreen}{yes} & on the training data only \\
C12 & direction: thallium stress test result: thallium stress test result other than a fixed defect on the thallium scan, a reversible defect on the thallium scan to a reversible defect on the thallium scan (whole model) \newline\textit{shape: categorical; highest to lowest: a reversible defect on the thallium scan, thallium stress test result other than a fixed defect on the thallium scan, a reversible defect on the thallium scan, a fixed defect on the thallium scan} & 30\% to 55\% & log-odds 1.03, OR 2.81; range across patients 1.03 to 1.03 & 137, 90 & train; test unclear (AUC gain 0.020 [-0.008, 0.054], n=61) & \textcolor{OliveGreen}{yes} & on the training data only \\
\bottomrule
\end{tabularx}}
\par{\tiny\noindent
\begin{tabularx}{\linewidth}{@{}l>{\raggedright\arraybackslash}X>{\raggedright\arraybackslash}p{0.11\linewidth}>{\raggedright\arraybackslash}p{0.16\linewidth}>{\raggedright\arraybackslash}p{0.05\linewidth}>{\raggedright\arraybackslash}p{0.15\linewidth}l>{\raggedright\arraybackslash}p{0.08\linewidth}@{}}
\toprule
Claim & What is checked & Risk, typical patient & Estimate (log-odds; OR); range across patients & n & Test check & Verified & Words \\
\midrule
C13 & direction: slope of the ST segment at peak exercise: an upsloping ST segment to a downsloping ST segment (whole model) \newline\textit{shape: categorical; highest to lowest: a flat ST segment, a downsloping ST segment, an upsloping ST segment} & 30\% to 54\% & log-odds 0.98, OR 2.66; range across patients 0.98 to 0.98 & 115, 17 (small n) & train; test disagrees (AUC gain -0.005 [-0.050, 0.037], n=61) & \textcolor{OliveGreen}{yes} & the model has little data on this; on the training data only \\
C14 & direction: chest pain type: no chest pain (asymptomatic) to non-anginal chest pain (direct contribution only, not through any part) \newline\textit{shape: categorical; highest to lowest: no chest pain (asymptomatic), atypical angina, non-anginal chest pain, typical angina} & 30\% to 7\% & log-odds -1.75, OR 0.17; range across patients -1.75 to -1.75 & 108, 70 & train; test agrees (AUC gain 0.038 [0.011, 0.072], n=61) & \textcolor{OliveGreen}{yes} & clearly \\
C15 & direction: thallium stress test result: thallium stress test result other than a fixed defect on the thallium scan, a reversible defect on the thallium scan to a fixed defect on the thallium scan (whole model) \newline\textit{shape: categorical; highest to lowest: a reversible defect on the thallium scan, thallium stress test result other than a fixed defect on the thallium scan, a reversible defect on the thallium scan, a fixed defect on the thallium scan} & 30\% to 12\% & log-odds -1.17, OR 0.31; range across patients -1.17 to -1.17 & 137, 15 (small n) & train; test unclear (AUC gain 0.020 [-0.008, 0.054], n=61) & \textcolor{OliveGreen}{yes} & the model has little data on this; on the training data only \\
C16 & direction: chest pain type: no chest pain (asymptomatic) to atypical angina (direct contribution only, not through any part) \newline\textit{shape: categorical; highest to lowest: no chest pain (asymptomatic), atypical angina, non-anginal chest pain, typical angina} & 30\% to 13\% & log-odds -1.03, OR 0.36; range across patients -1.03 to -1.03 & 108, 42 & train; test agrees (AUC gain 0.038 [0.011, 0.072], n=61) & \textcolor{OliveGreen}{yes} & clearly \\
C17 & direction: chest pain type: no chest pain (asymptomatic) to atypical angina (whole model) \newline\textit{shape: categorical; highest to lowest: no chest pain (asymptomatic), atypical angina, non-anginal chest pain, typical angina} & 30\% to 13\% & log-odds -1.03, OR 0.36; range across patients -1.03 to -1.03 & 108, 42 & train; test agrees (AUC gain 0.055 [0.014, 0.101], n=61) & \textcolor{OliveGreen}{yes} & clearly \\
C18 & direction: fasting blood sugar: fasting blood sugar of 120 mg/dL or less to fasting blood sugar above 120 mg/dL (whole model) \newline\textit{shape: categorical; highest to lowest: fasting blood sugar of 120 mg/dL or less, fasting blood sugar above 120 mg/dL} & 30\% to 19\% & log-odds -0.61, OR 0.54; range across patients -0.61 to -0.61 & 209, 33 & train; test unclear (AUC gain 0.001 [-0.012, 0.016], n=61) & \textcolor{OliveGreen}{yes} & on the training data only \\
C19 & direction: maximum heart rate: 118 beats per minute to 175 beats per minute (direct contribution only, not through any part) \newline\textit{shape: steady; 100\% of patients same; verified} & 36\% to 27\% & log-odds -0.45, OR 0.64; range across patients -0.45 to -0.45 & 242 & train; test unclear (AUC gain 0.005 [-0.004, 0.020], n=61) & \textcolor{OliveGreen}{yes} & on the training data only \\
C20 & direction: maximum heart rate: 118 beats per minute to 175 beats per minute (whole model) \newline\textit{shape: accelerating; 100\% of patients same; verified} & 36\% to 27\% & log-odds -0.70, OR 0.50; range across patients -1.33 to -0.44 & 242 & train; test unclear (AUC gain 0.012 [-0.002, 0.034], n=61) & \textcolor{OliveGreen}{yes} & on the training data only \\
C21 & direction: resting ECG: resting ECG other than signs of left ventricular hypertrophy on the resting ECG to signs of left ventricular hypertrophy on the resting ECG (whole model) \newline\textit{shape: categorical; highest to lowest: signs of left ventricular hypertrophy on the resting ECG, resting ECG other than signs of left ventricular hypertrophy on the resting ECG} & 30\% to 33\% & log-odds 0.11, OR 1.11; range across patients 0.11 to 0.11 & 123, 119 & train; test disagrees (AUC gain -0.002 [-0.010, 0.000], n=61) & \textcolor{OliveGreen}{yes} & slightly; on the training data only \\
C22 & magnitude: number of vessels coloured by fluoroscopy: 0 vessels to 2 vessels, typical patient \newline\textit{shape: steady (equal steps); 100\% of patients same; verified; risk saturates near 100\%} & 30\% to 96\% & log-odds 3.99, OR 54.1 & 242 & train; test unclear (AUC gain 0.039 [-0.034, 0.120], n=61) & \textcolor{OliveGreen}{yes} & on the training data only \\
C23 & magnitude: slope of the ST segment at peak exercise: an upsloping ST segment to a flat ST segment, typical patient \newline\textit{shape: categorical; highest to lowest: a flat ST segment, a downsloping ST segment, an upsloping ST segment} & 30\% to 80\% & log-odds 2.20, OR 9.04 & 115, 110 & train; test disagrees (AUC gain -0.005 [-0.050, 0.037], n=61) & \textcolor{OliveGreen}{yes} & on the training data only \\
C24 & magnitude: ST depression on exercise: 0.0 mm to 2.8 mm, typical patient \newline\textit{shape: steady; 100\% of patients same; verified} & 23\% to 53\% & log-odds 1.30, OR 3.68 & 242 & train; test unclear (AUC gain 0.011 [-0.01, 0.038], n=61) & \textcolor{OliveGreen}{yes} & on the training data only \\
\bottomrule
\end{tabularx}}
\par{\tiny\noindent
\begin{tabularx}{\linewidth}{@{}l>{\raggedright\arraybackslash}X>{\raggedright\arraybackslash}p{0.11\linewidth}>{\raggedright\arraybackslash}p{0.16\linewidth}>{\raggedright\arraybackslash}p{0.05\linewidth}>{\raggedright\arraybackslash}p{0.15\linewidth}l>{\raggedright\arraybackslash}p{0.08\linewidth}@{}}
\toprule
Claim & What is checked & Risk, typical patient & Estimate (log-odds; OR); range across patients & n & Test check & Verified & Words \\
\midrule
C25 & magnitude: chest pain type: no chest pain (asymptomatic) to typical angina, typical patient \newline\textit{shape: categorical; highest to lowest: no chest pain (asymptomatic), atypical angina, non-anginal chest pain, typical angina} & 30\% to 2\% & log-odds -3.11, OR 0.04 & 108, 22 & train; test agrees (AUC gain 0.055 [0.014, 0.101], n=61) & \textcolor{OliveGreen}{yes} & clearly \\
C26 & magnitude: serum cholesterol: 193 mg/dL (5.0 mmol/L) to 308 mg/dL (8.0 mmol/L), typical patient \newline\textit{shape: steady; 100\% of patients same; verified} & 21\% to 48\% & log-odds 1.26, OR 3.54 & 242 & train; test disagrees (AUC gain -0.012 [-0.044, 0.013], n=61) & \textcolor{OliveGreen}{yes} & on the training data only \\
C27 & null: model ignores age (exact rule) &  & max change 0.000 (max $|$change in log-odds$|$) & 242 & train & \textcolor{OliveGreen}{yes} &  \\
C28 & null: model ignores exercise-induced angina (exact rule) &  & max change 0.000 (max $|$change in log-odds$|$) & 242 & train & \textcolor{OliveGreen}{yes} &  \\
C29 & structure: made of nonanginal\_\allowbreak{}thalach\_\allowbreak{}gate,\allowbreak{} oldpeak\_\allowbreak{}gate\_\allowbreak{}without\_\allowbreak{}nonanginal,\allowbreak{} nonanginal\_\allowbreak{}low\_\allowbreak{}chol\_\allowbreak{}gate and inputs sex, cp, trestbps, chol, fbs, restecg, thalach, oldpeak, slope, ca, thal; added &  & log-odds n/a, OR n/a & 242 & train & \textcolor{OliveGreen}{yes} &  \\
C30 & direction: chest pain type: chest pain type other than non-anginal chest pain to non-anginal chest pain (through this part) \newline\textit{shape: categorical; highest to lowest: chest pain type other than non-anginal chest pain, non-anginal chest pain} & 30\% to 17\% & log-odds -0.69, OR 0.50; range across patients -1.07 to -0.19 & 172, 70 & train; test unclear (AUC gain 0.008 [-0.009, 0.028], n=61) & \textcolor{OliveGreen}{yes} & on the training data only \\
C31 & direction: maximum heart rate: 118 beats per minute to 175 beats per minute (through this part) \newline\textit{shape: hinge at 106 beats per minute; 29\% of patients same; fallback; risk saturates near 0\%} & 4\% to 2\% (typical patient this part acts on) & log-odds -0.25, OR 0.78; range across patients -0.88 to 0.00 & 242 & train; test unclear (AUC gain 0.006 [0.000, 0.020], n=61) & \textcolor{OliveGreen}{yes} & slightly; on the training data only \\
C32 & interaction: effect of maximum heart rate (118 beats per minute to 175 beats per minute) by chest pain type: non-anginal chest pain minus chest pain type other than non-anginal chest pain; stronger in non-anginal chest pain \newline\textit{shape: accelerating; 100\% of patients same; verified; risk saturates near 0\%} & non-anginal chest pain: 7\% to 2\%; chest pain type other than non-anginal chest pain: 36\% to 27\% & log-odds difference -0.89, OR 0.41 & 70, 172 & train; test agrees (AUC gain 0.008 [-0.009, 0.028], n=61); group difference on test -0.89 [-0.89, -0.88] & \textcolor{OliveGreen}{yes} & slightly \\
C33 & threshold: maximum heart rate: fold at about 184 beats per minute (starts) for patients with no chest pain symptoms (asymptomatic) \newline\textit{shape: step at 190 beats per minute; 100\% of patients same; verified; risk saturates near 100\%} & 99\% to 99\% (typical patient with a fold (patients with no chest pain symptoms (asymptomatic))) & fold 184 beats per minute (IQR 184 to 184); past it log-odds 0.00, OR 1.00; range across patients 0.00 to 0.00 & 108, 3 (small n) & train; test unclear (AUC gain 0.006 [0.000, 0.020], n=61) & \textcolor{BrickRed}{no} & slightly; the model has little data on this; on the training data only \\
C34 & subgroup: who this part acts on: 76 of 242 patients & 5\% to 2\% (typical patient this part acts on (risk without the part, then with it)) & share 0.31; log-odds -0.71, OR 0.49; range across patients -1.06 to -0.14 & 76, 166 & train; test unclear (AUC gain 0.013 [-0.001, 0.038], n=61) & \textcolor{OliveGreen}{yes} & slightly; on the training data only \\
C35 & structure: its formula reproduces it on 242 training patients (largest difference 4.44089e-16) &  & log-odds 4e-16, OR n/a & 242 & train & \textcolor{OliveGreen}{yes} &  \\
C36 & direction: chest pain type: chest pain type other than non-anginal chest pain to non-anginal chest pain (through this part) \newline\textit{shape: categorical; highest to lowest: chest pain type other than non-anginal chest pain, non-anginal chest pain} & 30\% to 29\% & log-odds -0.08, OR 0.92; range across patients -0.16 to 0.00 & 172, 70 & train; test unclear (AUC gain 0.002 [0.000, 0.009], n=61) & \textcolor{OliveGreen}{yes} & slightly; on the training data only \\
\bottomrule
\end{tabularx}}
\par{\tiny\noindent
\begin{tabularx}{\linewidth}{@{}l>{\raggedright\arraybackslash}X>{\raggedright\arraybackslash}p{0.11\linewidth}>{\raggedright\arraybackslash}p{0.16\linewidth}>{\raggedright\arraybackslash}p{0.05\linewidth}>{\raggedright\arraybackslash}p{0.15\linewidth}l>{\raggedright\arraybackslash}p{0.08\linewidth}@{}}
\toprule
Claim & What is checked & Risk, typical patient & Estimate (log-odds; OR); range across patients & n & Test check & Verified & Words \\
\midrule
C37 & direction: serum cholesterol: 193 mg/dL (5.0 mmol/L) to 308 mg/dL (8.0 mmol/L) (through this part) \newline\textit{shape: hinge at 282 mg/dL (7.3 mmol/L); 100\% of patients same; verified; risk saturates near 0\%; one point with C39 (median fold 282 mg/dL (7.3 mmol/L), typical patient 283 mg/dL (7.3 mmol/L), 10th-90th percentile 282 mg/dL (7.3 mmol/L) to 282 mg/dL (7.3 mmol/L))} & 2\% to 3\% (typical patient this part acts on) & log-odds 0.05, OR 1.05; range across patients 0.00 to 0.16 & 242 & train; test disagrees (AUC gain -0.001 [-0.006, 0.000], n=61) & \textcolor{OliveGreen}{yes} & slightly; on the training data only \\
C38 & interaction: effect of serum cholesterol (193 mg/dL (5.0 mmol/L) to 308 mg/dL (8.0 mmol/L)) by chest pain type: non-anginal chest pain minus chest pain type other than non-anginal chest pain; stronger in non-anginal chest pain \newline\textit{shape: steady; 100\% of patients same; fallback; risk saturates near 0\%} & non-anginal chest pain: 2\% to 7\%; chest pain type other than non-anginal chest pain: 21\% to 48\% & log-odds difference 0.16, OR 1.18 & 70, 172 & train; test agrees (AUC gain 0.002 [0.000, 0.009], n=61); group difference on test 0.16 [0.16, 0.16] & \textcolor{OliveGreen}{yes} & clearly \\
C39 & threshold: serum cholesterol: fold at about 282 mg/dL (7.3 mmol/L) (stops) for patients with non-anginal chest pain \newline\textit{shape: hinge at 282 mg/dL (7.3 mmol/L); 100\% of patients same; verified; risk saturates near 0\%; one point with C39 (median fold 282 mg/dL (7.3 mmol/L), typical patient 283 mg/dL (7.3 mmol/L), 10th-90th percentile 282 mg/dL (7.3 mmol/L) to 282 mg/dL (7.3 mmol/L))} & 2\% to 3\% (typical patient with a fold) & fold 282 mg/dL (IQR 282 to 282); past it log-odds 0.19, OR 1.21; range across patients 0.19 to 0.19 & 70, 61 & train; test disagrees (AUC gain -0.001 [-0.006, 0.000], n=61) & \textcolor{OliveGreen}{yes} & slightly; on the training data only \\
C40 & subgroup: who this part acts on: 61 of 242 patients & 3\% to 2\% (typical patient this part acts on (risk without the part, then with it)) & share 0.25; log-odds -0.10, OR 0.91; range across patients -0.18 to -0.02 & 61, 181 & train; test unclear (AUC gain 0.000 [0.000, 0.000], n=61) & \textcolor{OliveGreen}{yes} & slightly; on the training data only \\
C41 & structure: its formula reproduces it on 242 training patients (largest difference 2.22045e-16) &  & log-odds 2e-16, OR n/a & 242 & train & \textcolor{OliveGreen}{yes} &  \\
C42 & direction: ST depression on exercise: 0.0 mm to 2.8 mm (through this part) \newline\textit{shape: flat; 100\% of patients same; verified} & 30\% to 30\% & log-odds -0.01, OR 0.99; range across patients -0.02 to 0.00 & 242 & train; test unclear (AUC gain 0.000 [0.000, 0.000], n=61) & \textcolor{OliveGreen}{yes} & slightly; on the training data only \\
C43 & direction: maximum heart rate: 118 beats per minute to 175 beats per minute (through this part) \newline\textit{shape: flat; 100\% of patients same; verified} & 30\% to 30\% & log-odds 0.004, OR 1.00; range across patients 0.00 to 0.01 & 242 & train; test unclear (AUC gain 0.000 [0.000, 0.000], n=61) & \textcolor{OliveGreen}{yes} & slightly; on the training data only \\
C44 & direction: chest pain type: chest pain type other than non-anginal chest pain to non-anginal chest pain (through this part) \newline\textit{shape: categorical; highest to lowest: non-anginal chest pain, chest pain type other than non-anginal chest pain} & $>$99\% to $>$99\% (typical patient this part acts on) & log-odds 0.005, OR 1.01; range across patients 0.00 to 0.02 & 172, 70 & train; test unclear (AUC gain 0.000 [0.000, 0.000], n=61) & \textcolor{OliveGreen}{yes} & slightly; on the training data only \\
C45 & interaction: effect of ST depression on exercise (0.0 mm to 2.8 mm) by chest pain type: non-anginal chest pain minus chest pain type other than non-anginal chest pain; stronger in non-anginal chest pain \newline\textit{shape: steady; 100\% of patients same; verified; risk saturates near 0\%} & non-anginal chest pain: 2\% to 8\%; chest pain type other than non-anginal chest pain: 23\% to 53\% & log-odds difference 0.02, OR 1.02 & 70, 172 & train; test agrees (AUC gain 0.000 [0.000, 0.000], n=61); group difference on test 0.02 [0.02, 0.02] & \textcolor{BrickRed}{no} & clearly \\
C46 & interaction: effect of ST depression on exercise (0.0 mm to 2.8 mm) by maximum heart rate: maximum heart rate above 152 beats per minute minus maximum heart rate at or below 152 beats per minute; stronger in maximum heart rate above 152 beats per minute \newline\textit{shape: steady; 100\% of patients same; verified} & maximum heart rate above 152 beats per minute: 21\% to 50\%; maximum heart rate at or below 152 beats per minute: 26\% to 56\% & log-odds difference 0.007, OR 1.01 & 121, 121 & train; test agrees (AUC gain 0.000 [0.000, 0.000], n=61); group difference on test 0.008 [0.004, 0.01] & \textcolor{BrickRed}{no} & slightly \\
C47 & threshold: ST depression on exercise: fold at about 1.0 mm (starts) for patients with no chest pain symptoms (asymptomatic) \newline\textit{shape: flat; 100\% of patients same; verified; risk saturates near 100\%} & 99\% to 99\% (typical patient with a fold (patients with no chest pain symptoms (asymptomatic))) & fold 1.01 mm (IQR 0.67 to 1.31); past it log-odds -0.03, OR 0.97; range across patients -0.03 to -0.02 & 106, 58 & train; test unclear (AUC gain 0.000 [0.000, 0.000], n=61) & \textcolor{BrickRed}{no} & slightly; on the training data only \\
C48 & threshold: maximum heart rate: fold at about 90 beats per minute (stops) for patients with atypical angina \newline\textit{shape: flat; 100\% of patients same; verified} & 9\% to 9\% (typical patient with a fold (patients with atypical angina)) & fold 90.0 beats per minute (IQR 90.0 to 104); past it log-odds 0.002, OR 1.00; range across patients 0.00 to 0.008 & 42, 3 (small n) & train; test unclear (AUC gain 0.000 [0.000, 0.000], n=61) & \textcolor{BrickRed}{no} & slightly; the model has little data on this; on the training data only \\
\bottomrule
\end{tabularx}}
\par{\tiny\noindent
\begin{tabularx}{\linewidth}{@{}l>{\raggedright\arraybackslash}X>{\raggedright\arraybackslash}p{0.11\linewidth}>{\raggedright\arraybackslash}p{0.16\linewidth}>{\raggedright\arraybackslash}p{0.05\linewidth}>{\raggedright\arraybackslash}p{0.15\linewidth}l>{\raggedright\arraybackslash}p{0.08\linewidth}@{}}
\toprule
Claim & What is checked & Risk, typical patient & Estimate (log-odds; OR); range across patients & n & Test check & Verified & Words \\
\midrule
C49 & subgroup: who this part acts on: 73 of 242 patients & $>$99\% to $>$99\% (typical patient this part acts on (risk without the part, then with it)) & share 0.30; log-odds -0.02, OR 0.98; range across patients -0.03 to -0.002 & 73, 169 & train; test unclear (AUC gain 0.000 [0.000, 0.000], n=61) & \textcolor{BrickRed}{no} & slightly; on the training data only \\
C50 & null: this part barely changes predictions (switched off) &  & max change 0.005 (max $|$change in predicted risk$|$) & 242 & train & \textcolor{OliveGreen}{yes} &  \\
C51 & structure: its formula reproduces it on 242 training patients (largest difference 2.22045e-16) &  & log-odds 2e-16, OR n/a & 242 & train & \textcolor{OliveGreen}{yes} &  \\
C52 & boundary: boundary &  & log-odds 1.00, OR n/a & 242 & train & \textcolor{OliveGreen}{yes} &  \\
\bottomrule
\end{tabularx}}
\paragraph{Not supported by the model.}
\begin{itemize}\setlength\itemsep{0pt}\small
\item C33: threshold claim about maximum heart rate reached during the exercise test in part ann\_\allowbreak{}30: not supported by the model (the effect past the fold is too small or not consistent across patients).
\item C45: interaction claim about ST depression during exercise (compared with rest) and chest pain type in part ann\_\allowbreak{}31: not supported by the model (the group difference is below 0.1 log-odds).
\item C46: interaction claim about ST depression during exercise (compared with rest) and maximum heart rate reached during the exercise test in part ann\_\allowbreak{}31: not supported by the model (the group difference is below 0.1 log-odds).
\item C47: threshold claim about ST depression during exercise (compared with rest) in part ann\_\allowbreak{}31: not supported by the model (the fold changes the curve by only 0.011 log-odds compared with a straight line (needed 0.1); the effect past the fold is too small or not consistent across patients).
\item C48: threshold claim about maximum heart rate reached during the exercise test in part ann\_\allowbreak{}31: not supported by the model (the fold changes the curve by only 0.022 log-odds compared with a straight line (needed 0.1); the effect past the fold is too small or not consistent across patients).
\item C49: subgroup claim about chest pain type and maximum heart rate reached during the exercise test in part ann\_\allowbreak{}31: not supported by the model (its effect is negligible).
\end{itemize}

\subsection{Coverage}\label{app:heart-s2-c3:coverage}
Coverage is computed by the current code (the function behind the \texttt{get\_coverage} tool) under these conventions. \emph{Candidate nodes} are all nodes of the explained model state except the output nodes. A candidate node is \emph{covered} when some annotation of the hierarchy (a leaf, or a composition together with its descendants) contains it and all its enabled outgoing edges. \emph{Exact Null} parts, which have no effect on any output for any input, count as covered: nodes with no path to an output over enabled, non-zero-weight edges (disconnected inputs and dead ends), and edges that have zero weight or touch such a node. Compositional coverage is the fraction of required compositions present: every annotation except a single global root containing every output must have a parent; a parentless, childless annotation made only of Null nodes is exempt.

\begin{itemize}
\item \textbf{Null, disconnected input (9):} \texttt{-1} (age), \texttt{-10} (fbs:0), \texttt{-12} (restecg:0), \texttt{-13} (restecg:1), \texttt{-16} (exang:0), \texttt{-17} (exang:1), \texttt{-21} (slope:3), \texttt{-23} (thal:3), \texttt{-24} (thal:4).
\item \textbf{Uncovered nodes:} none.
\item \textbf{Uncovered edges:} none.
\item \textbf{Missing compositions:} none.
\end{itemize}

\subsection{Annotation hierarchy}\label{app:heart-s2-c3:hierarchy}
The hierarchy has 1 root; the global root is heart\_\allowbreak{}disease\_\allowbreak{}logit. Tags (A1, A2, \ldots) are used in the figures.
\begingroup\parindent=0pt\parskip=1pt
\hspace*{0.0em}$\triangleright$~A1: heart\_\allowbreak{}disease\_\allowbreak{}logit (composition, \cref{app:heart-s2-c3:ann-heart-disease-logit})\par
\hspace*{1.5em}$\triangleright$~A2: nonanginal\_\allowbreak{}thalach\_\allowbreak{}gate (leaf, \cref{app:heart-s2-c3:ann-nonanginal-thalach-gate})\par
\hspace*{1.5em}$\triangleright$~A3: oldpeak\_\allowbreak{}gate\_\allowbreak{}without\_\allowbreak{}nonanginal (leaf, \cref{app:heart-s2-c3:ann-oldpeak-gate-without-nonanginal})\par
\hspace*{1.5em}$\triangleright$~A4: nonanginal\_\allowbreak{}low\_\allowbreak{}chol\_\allowbreak{}gate (leaf, \cref{app:heart-s2-c3:ann-nonanginal-low-chol-gate})\par
\endgroup

\subsection{Network figures}\label{app:heart-s2-c3:figures}
\begin{figure}[htbp]
\centering
\input{\explaneatcertdir{heart-s2-c3}heart-s2-c3-fig-full.tex}
\caption{heart-s2-c3: full view of the explained model state.}
\label{app:heart-s2-c3:fig-full}
\end{figure}
\begin{figure}[htbp]
\centering
\input{\explaneatcertdir{heart-s2-c3}heart-s2-c3-fig-level-0.tex}
\caption{heart-s2-c3: collapsed view at hierarchy level 0 (fully recomposed).}
\label{app:heart-s2-c3:fig-level-0}
\end{figure}
\begin{figure}[htbp]
\centering
\input{\explaneatcertdir{heart-s2-c3}heart-s2-c3-fig-level-1.tex}
\caption{heart-s2-c3: collapsed view at hierarchy level 1.}
\label{app:heart-s2-c3:fig-level-1}
\end{figure}

\subsection{Refactoring operations}\label{app:heart-s2-c3:operations}
The explained model state is the phenotype after these operations, replayed in order. All are function-preserving; the largest change in the model output over the 242 training rows is $2.22\times 10^{-16}$ (tolerance $1.00\times 10^{-9}$).
\begin{itemize}
\item \#26: split \texttt{-6} into \texttt{-6\_\allowbreak{}a}, \texttt{-6\_\allowbreak{}b}, \texttt{-6\_\allowbreak{}c}.
\item \#27: split \texttt{-15} into \texttt{-15\_\allowbreak{}a}, \texttt{-15\_\allowbreak{}b}, \texttt{-15\_\allowbreak{}c}.
\item \#28: split \texttt{-18} into \texttt{-18\_\allowbreak{}a}, \texttt{-18\_\allowbreak{}b}.
\item \#29: split \texttt{-9} into \texttt{-9\_\allowbreak{}a}, \texttt{-9\_\allowbreak{}b}.
\end{itemize}

\subsection{Annotations}\label{app:heart-s2-c3:annotations}
\subsubsection{\texorpdfstring{A1: heart\_\allowbreak{}disease\_\allowbreak{}logit}{A1: heart disease logit}}\label{app:heart-s2-c3:ann-heart-disease-logit}
% annotation id ann_33, operation seq 33
\textbf{Kind:} Composition annotation (depth 0); children A2 (\cref{app:heart-s2-c3:ann-nonanginal-thalach-gate}), A3 (\cref{app:heart-s2-c3:ann-oldpeak-gate-without-nonanginal}), A4 (\cref{app:heart-s2-c3:ann-nonanginal-low-chol-gate}).

\begin{itemize}
\item \textbf{Entry nodes:} \texttt{-2} (sex:0), \texttt{-3} (sex:1), \texttt{-4} (cp:1), \texttt{-5} (cp:2), \texttt{-7} (cp:4), \texttt{-8} (trestbps), \texttt{-9\_\allowbreak{}a} (chol\_\allowbreak{}a), \texttt{-11} (fbs:1), \texttt{-14} (restecg:2), \texttt{-15\_\allowbreak{}a} (thalach\_\allowbreak{}a), \texttt{-18\_\allowbreak{}a} (oldpeak\_\allowbreak{}a), \texttt{-19} (slope:1), \texttt{-20} (slope:2), \texttt{-22} (ca), \texttt{-25} (thal:6), \texttt{-26} (thal:7), \texttt{-1} (age), \texttt{-10} (fbs:0), \texttt{-12} (restecg:0), \texttt{-13} (restecg:1), \texttt{-16} (exang:0), \texttt{-17} (exang:1), \texttt{-21} (slope:3), \texttt{-23} (thal:3), \texttt{-24} (thal:4), \texttt{-6\_\allowbreak{}a} (cp:3\_\allowbreak{}a), \texttt{-15\_\allowbreak{}b} (thalach\_\allowbreak{}b), \texttt{-6\_\allowbreak{}b} (cp:3\_\allowbreak{}b), \texttt{-15\_\allowbreak{}c} (thalach\_\allowbreak{}c), \texttt{-18\_\allowbreak{}b} (oldpeak\_\allowbreak{}b), \texttt{-6\_\allowbreak{}c} (cp:3\_\allowbreak{}c), \texttt{-9\_\allowbreak{}b} (chol\_\allowbreak{}b).
\item \textbf{Exit nodes:} \texttt{0}.
\item \textbf{Internal nodes:} none.
\item \textbf{Edges (weight):} \texttt{-2}$\to$\texttt{0}~(-0.76), \texttt{-3}$\to$\texttt{0}~(0.347), \texttt{-4}$\to$\texttt{0}~(-0.393), \texttt{-5}$\to$\texttt{0}~(0.273), \texttt{-7}$\to$\texttt{0}~(0.869), \texttt{-8}$\to$\texttt{0}~(0.493), \texttt{-9\_\allowbreak{}a}$\to$\texttt{0}~(0.529), \texttt{-11}$\to$\texttt{0}~(-0.209), \texttt{-14}$\to$\texttt{0}~(0.053), \texttt{-15\_\allowbreak{}a}$\to$\texttt{0}~(-0.176), \texttt{-18\_\allowbreak{}a}$\to$\texttt{0}~(0.559), \texttt{-19}$\to$\texttt{0}~(-0.488), \texttt{-20}$\to$\texttt{0}~(0.61), \texttt{-22}$\to$\texttt{0}~(1.78), \texttt{-25}$\to$\texttt{0}~(-0.282), \texttt{-26}$\to$\texttt{0}~(0.5), \texttt{2472}$\to$\texttt{0}~(-0.511), \texttt{2596}$\to$\texttt{0}~(-0.0184), \texttt{2863}$\to$\texttt{0}~(-0.235).
\item \textbf{With its descendants:} 36 nodes.
\item \textbf{Covers} (coverage, own and as part of the hierarchy): 35 candidate nodes.
\end{itemize}

\textbf{Hypothesis.} Root: p = sigmoid(z), z = 0.323 + LINEAR CORE - 0.511*\allowbreak{}nonanginal\_\allowbreak{}thalach\_\allowbreak{}gate - 0.018*\allowbreak{}oldpeak\_\allowbreak{}gate\_\allowbreak{}without\_\allowbreak{}nonanginal - 0.235*\allowbreak{}nonanginal\_\allowbreak{}low\_\allowbreak{}chol\_\allowbreak{}gate. LINEAR CORE (16 direct input-$>$output edges, z-scored inputs): ca +1.779, cp:4 +0.869, sex:0 -0.760, slope:2 +0.610, oldpeak +0.559, chol +0.529, thal:7 +0.500, trestbps +0.493, slope:1 -0.488, cp:1 -0.393, sex:1 +0.347, thal:6 -0.282, cp:2 +0.273, fbs:1 -0.209, thalach -0.176, restecg:2 +0.053. Nine inputs are unused (age, fbs:0, restecg:0, restecg:1, exang:0, exang:1, slope:3, thal:3, thal:4; zero-effect entries of this root). Non-anginal pain (cp:3) has no direct edge; it acts only through the three children, all of which are switched by it. COMPOSITION: purely additive in the logit. The children share no hidden nodes and are coupled only through shared inputs (cp:3 feeds all three; thalach, oldpeak and chol also have direct edges). (1) nonanginal\_\allowbreak{}thalach\_\allowbreak{}gate: open for 97\% of non-anginal-pain patients, lowers their risk, more with higher max heart rate; (2) nonanginal\_\allowbreak{}low\_\allowbreak{}chol\_\allowbreak{}gate: open for 87\% of non-anginal-pain patients (those with z-chol $<$ 0.76), lowers their risk more the lower the cholesterol; (3) oldpeak\_\allowbreak{}gate\_\allowbreak{}without\_\allowbreak{}nonanginal: a near-null oldpeak hinge for the other patients (readout -0.018). Together (1)+(2) give non-anginal pain an average logit shift of -0.85 and make thalach and chol slopes depend on chest-pain type. EVIDENCE (train, n=242): full model AUC 0.9381, acc 0.888 (compute\_\allowbreak{}performance)\allowbreak{}. Linear core alone AUC 0.934; logit SD 3.67 (core) vs 0.42 (all nonlinear readouts). Zero-ablating one child at a time: thalach gate 0.934, oldpeak gate 0.938, low-chol gate 0.938. Replacing all hidden readouts with one least-squares linear fit in their four inputs gives 0.938; an L2 logistic regression on all 26 inputs gives 0.936. Mean $|$SHAP$|$ (train, all 242 rows): ca 0.161, cp:4 0.099, sex:0 0.073, slope:2 0.066, thal:7 0.054, slope:1 0.053, oldpeak 0.047, chol 0.041, trestbps 0.040, cp:3 (thalach gate copy) 0.029; unused inputs about 0.001 (estimator noise). Interpretation: logistic regression on 16 inputs plus a non-anginal-pain block whose nonlinearity is small (a linear refit loses nothing on train AUC). Counter-directional signs reported, not explained: cp:2 +0.27 and fbs:1 -0.21 in the core, the low-cholesterol protection for non-anginal pain. Compared with the 10-seed heart-s2 record (key 9083, gen 38, a different cluster of the same run): that genome has an age/\allowbreak{}heart-rate-reserve gate and oldpeak hinges and uses age; this one ignores age entirely and puts its only nonlinear effect on non-anginal chest pain.

\textbf{Closed form} (children as named functions):
{\small\begingroup\allowdisplaybreaks
\begin{align*}
y_{0} &= \sigma\bigl(z_{0}\bigr) \\
z_{0} &= 0.323 + 0.559 x_{oldpeak\_a} + 0.529 x_{chol\_a} + 0.053 x_{restecg:2} + 1.78 x_{ca} \\
&\qquad + 0.347 x_{sex:1} + 0.273 x_{cp:2} + 0.61 x_{slope:2} + 0.869 x_{cp:4} + 0.493 x_{trestbps} \\
&\qquad + 0.5 x_{thal:7} - 0.176 x_{thalach\_a} \\
&\qquad - 0.235 \operatorname{nonanginal\_low\_chol\_gate}\bigl(x_{cp:3\_c}, x_{chol\_b}\bigr) \\
&\qquad - 0.488 x_{slope:1} - 0.511 \operatorname{nonanginal\_thalach\_gate}\bigl(x_{cp:3\_a}, \\
&\qquad x_{thalach\_b}\bigr) - 0.393 x_{cp:1} - 0.209 x_{fbs:1} - 0.282 x_{thal:6} \\
&\qquad - 0.0184 \operatorname{oldpeak\_gate\_without\_nonanginal}\bigl(x_{cp:3\_b}, x_{thalach\_c}, \\
&\qquad x_{oldpeak\_b}\bigr) - 0.76 x_{sex:0}
\end{align*}\endgroup}

Expanded over the children:
{\small\begingroup\allowdisplaybreaks
\begin{align*}
y_{0} &= \sigma\bigl(z_{0}\bigr) \\
z_{0} &= 0.323 + 0.559 x_{oldpeak\_a} + 0.529 x_{chol\_a} + 0.053 x_{restecg:2} + 1.78 x_{ca} \\
&\qquad + 0.347 x_{sex:1} + 0.273 x_{cp:2} + 0.61 x_{slope:2} + 0.869 x_{cp:4} + 0.493 x_{trestbps} \\
&\qquad + 0.5 x_{thal:7} - 0.176 x_{thalach\_a} - 0.235 \operatorname{relu}\bigl(-1.25 \\
&\qquad + 0.983 x_{cp:3\_c} - 0.376 x_{chol\_b}\bigr) - 0.488 x_{slope:1} \\
&\qquad - 0.511 \operatorname{relu}\bigl(-0.348 + 1.08 x_{cp:3\_a} + 0.678 x_{thalach\_b}\bigr) \\
&\qquad - 0.393 x_{cp:1} - 0.209 x_{fbs:1} - 0.282 x_{thal:6} - 0.0184 \operatorname{relu}\bigl(-0.688 \\
&\qquad + 0.673 x_{oldpeak\_b} - 0.241 x_{thalach\_c} - 0.985 x_{cp:3\_b}\bigr) - 0.76 x_{sex:0}
\end{align*}\endgroup}

where $y_{0}$ is node \texttt{0}; $z_{i}$, where shown, is the pre-activation of a sigmoid output $y_{i}$; weights and biases are shown to 3 significant figures.
 On the training split (n = 242), the displayed closed form differs from the forward pass by at most 0.00141, and the exact function by at most $2.22\times 10^{-16}$.

\textbf{Evidence} (2 items; full payloads in the JSON record):
\begin{itemize}
\item \textsc{analytical}, \texttt{analytical\_\allowbreak{}formula}; data split: n/a (analytical). Root closed form (get\_\allowbreak{}formula force=true); output is sigmoid.
\item \textsc{empirical}, \texttt{empirical\_\allowbreak{}viz}; data split: train. Train-split performance, ablations, linear replacements and SHAP.
\end{itemize}

\textbf{Authorship:} drafted by AI-assisted draft on 2026-11-10; verified by nobody (unverified).

\subsubsection{\texorpdfstring{A2: nonanginal\_\allowbreak{}thalach\_\allowbreak{}gate}{A2: nonanginal thalach gate}}\label{app:heart-s2-c3:ann-nonanginal-thalach-gate}
% annotation id ann_30, operation seq 30
\textbf{Kind:} Leaf annotation (depth 1); parent A1 (\cref{app:heart-s2-c3:ann-heart-disease-logit}).

\begin{itemize}
\item \textbf{Entry nodes:} \texttt{-6\_\allowbreak{}a} (cp:3\_\allowbreak{}a), \texttt{-15\_\allowbreak{}b} (thalach\_\allowbreak{}b).
\item \textbf{Exit nodes:} \texttt{2472}.
\item \textbf{Internal nodes:} none.
\item \textbf{Edges (weight):} \texttt{-6\_\allowbreak{}a}$\to$\texttt{2472}~(1.08), \texttt{-15\_\allowbreak{}b}$\to$\texttt{2472}~(0.678).
\item \textbf{Covers} (coverage, own and as part of the hierarchy): 2 candidate nodes.
\end{itemize}

\textbf{Hypothesis.} Non-anginal-pain gate (protective). 2472 = ReLU(-0.348 + 1.076*cp:3 + 0.678*thalach), read out -0.511 into the logit. Hypothesis: this gate is how the model gives non-anginal chest pain (cp:3, which has no direct edge) a lower risk: it is open for nearly every non-anginal-pain patient and almost never otherwise, and once open a higher achieved max heart rate lowers the risk further (on top of thalach's small direct -0.176 edge). EVIDENCE (train, n=242): active for 97\% of the 70 non-anginal-pain patients (closed only for z-thalach $<$ -1.97) and 5\% of the 172 others (only z-thalach $>$ 1.53). Disease rate 20\% when active vs 55\% when inactive. For a non-anginal-pain patient it adds -0.684 - 0.346*z-thalach to the logit (mean -0.76). Mean $|$sensitivity$|$ cp:3 0.34 $>$ thalach 0.21. Zero-ablating its readout lowers train AUC 0.938 -$>$ 0.934, the only nonlinear part that matters; a linear refit in (cp:3, thalach) recovers 0.938 (R2 0.87), so the gate acts almost as a cp:3 main effect plus a cp:3-specific thalach slope. Supported.

\textbf{Closed form:}
{\small\begingroup\allowdisplaybreaks
\begin{align*}
y_{0} &= \operatorname{relu}\bigl(-0.348 + 1.08 x_{cp:3\_a} + 0.678 x_{thalach\_b}\bigr)
\end{align*}\endgroup}

where $y_{0}$ is node \texttt{2472}; $z_{i}$, where shown, is the pre-activation of a sigmoid output $y_{i}$; weights and biases are shown to 3 significant figures.
 On the training split (n = 242), the displayed closed form differs from the forward pass by at most 0.00675, and the exact function by at most 0.

\textbf{Evidence} (2 items; full payloads in the JSON record):
\begin{itemize}
\item \textsc{analytical}, \texttt{analytical\_\allowbreak{}formula}; data split: n/a (analytical). Closed form from get\_\allowbreak{}formula with the on-data reading of the gate.
\item \textsc{empirical}, \texttt{empirical\_\allowbreak{}viz} (sensitivity (max\_\allowbreak{}samples 100000, sample\_\allowbreak{}fraction 1.0)); data split: train. Train-split sensitivity, activity, outcome rates and ablation.
\end{itemize}

\textbf{Authorship:} drafted by AI-assisted draft on 2026-11-10; verified by nobody (unverified).

\subsubsection{\texorpdfstring{A3: oldpeak\_\allowbreak{}gate\_\allowbreak{}without\_\allowbreak{}nonanginal}{A3: oldpeak gate without nonanginal}}\label{app:heart-s2-c3:ann-oldpeak-gate-without-nonanginal}
% annotation id ann_31, operation seq 31
\textbf{Kind:} Leaf annotation (depth 1); parent A1 (\cref{app:heart-s2-c3:ann-heart-disease-logit}).

\begin{itemize}
\item \textbf{Entry nodes:} \texttt{-6\_\allowbreak{}b} (cp:3\_\allowbreak{}b), \texttt{-15\_\allowbreak{}c} (thalach\_\allowbreak{}c), \texttt{-18\_\allowbreak{}b} (oldpeak\_\allowbreak{}b).
\item \textbf{Exit nodes:} \texttt{2596}.
\item \textbf{Internal nodes:} none.
\item \textbf{Edges (weight):} \texttt{-6\_\allowbreak{}b}$\to$\texttt{2596}~(-0.985), \texttt{-15\_\allowbreak{}c}$\to$\texttt{2596}~(-0.241), \texttt{-18\_\allowbreak{}b}$\to$\texttt{2596}~(0.673).
\item \textbf{Covers} (coverage, own and as part of the hierarchy): 3 candidate nodes.
\end{itemize}

\textbf{Hypothesis.} ST-depression gate for patients without non-anginal pain (negligible). 2596 = ReLU(-0.688 - 0.985*cp:3 - 0.241*thalach + 0.673*oldpeak), read out -0.018. Hypothesis: a hinge on ST depression (and lower max heart rate) for patients without non-anginal pain, whose readout is too small to change predictions. EVIDENCE (train, n=242): active for 42\% of patients without non-anginal pain and never with it; disease rate 80\% when active vs 29\% when inactive, so it tracks risk, but its readout has SD 0.009 logit units against 3.67 for the linear core, and removing it changes train AUC by +0.0001. Its sign (slightly lowering the risk with more ST depression) runs against oldpeak's direct +0.559 edge, which dominates. Supported as a near-null part; a vestige rather than a mechanism.

\textbf{Closed form:}
{\small\begingroup\allowdisplaybreaks
\begin{align*}
y_{0} &= \operatorname{relu}\bigl(-0.688 + 0.673 x_{oldpeak\_b} - 0.241 x_{thalach\_c} \\
&\qquad - 0.985 x_{cp:3\_b}\bigr)
\end{align*}\endgroup}

where $y_{0}$ is node \texttt{2596}; $z_{i}$, where shown, is the pre-activation of a sigmoid output $y_{i}$; weights and biases are shown to 3 significant figures.
 On the training split (n = 242), the displayed closed form differs from the forward pass by at most 0.00186, and the exact function by at most 0.

\textbf{Evidence} (2 items; full payloads in the JSON record):
\begin{itemize}
\item \textsc{analytical}, \texttt{analytical\_\allowbreak{}formula}; data split: n/a (analytical). Closed form from get\_\allowbreak{}formula with the on-data reading of the gate.
\item \textsc{empirical}, \texttt{empirical\_\allowbreak{}viz} (sensitivity (max\_\allowbreak{}samples 100000, sample\_\allowbreak{}fraction 1.0)); data split: train. Train-split sensitivity, activity, outcome rates and ablation.
\end{itemize}

\textbf{Authorship:} drafted by AI-assisted draft on 2026-11-10; verified by nobody (unverified).

\subsubsection{\texorpdfstring{A4: nonanginal\_\allowbreak{}low\_\allowbreak{}chol\_\allowbreak{}gate}{A4: nonanginal low chol gate}}\label{app:heart-s2-c3:ann-nonanginal-low-chol-gate}
% annotation id ann_32, operation seq 32
\textbf{Kind:} Leaf annotation (depth 1); parent A1 (\cref{app:heart-s2-c3:ann-heart-disease-logit}).

\begin{itemize}
\item \textbf{Entry nodes:} \texttt{-6\_\allowbreak{}c} (cp:3\_\allowbreak{}c), \texttt{-9\_\allowbreak{}b} (chol\_\allowbreak{}b).
\item \textbf{Exit nodes:} \texttt{2863}.
\item \textbf{Internal nodes:} none.
\item \textbf{Edges (weight):} \texttt{-6\_\allowbreak{}c}$\to$\texttt{2863}~(0.983), \texttt{-9\_\allowbreak{}b}$\to$\texttt{2863}~(-0.376).
\item \textbf{Covers} (coverage, own and as part of the hierarchy): 2 candidate nodes.
\end{itemize}

\textbf{Hypothesis.} Non-anginal-pain, low-cholesterol gate (protective, small). 2863 = ReLU(-1.255 + 0.983*cp:3 - 0.376*chol), read out -0.235. Hypothesis: open only for non-anginal-pain patients whose cholesterol is not high (z-chol below 0.76); it lowers their risk, by more the lower the cholesterol, so for these patients cholesterol's direct +0.529 edge is partly undone. EVIDENCE (train, n=242): active for 87\% of the 70 non-anginal-pain patients and none of the others. Disease rate 21\% when active vs 52\% when inactive. Readout up to -0.29 logit units (SD 0.053). Mean $|$sensitivity$|$ cp:3 0.25 $>$ chol 0.09. Zero-ablating it leaves train AUC unchanged (0.9381 -$>$ 0.9382): it shifts the level of the non-anginal-pain group rather than ranking patients. Supported but small.

\textbf{Closed form:}
{\small\begingroup\allowdisplaybreaks
\begin{align*}
y_{0} &= \operatorname{relu}\bigl(-1.25 + 0.983 x_{cp:3\_c} - 0.376 x_{chol\_b}\bigr)
\end{align*}\endgroup}

where $y_{0}$ is node \texttt{2863}; $z_{i}$, where shown, is the pre-activation of a sigmoid output $y_{i}$; weights and biases are shown to 3 significant figures.
 On the training split (n = 242), the displayed closed form differs from the forward pass by at most 0.00471, and the exact function by at most 0.

\textbf{Evidence} (2 items; full payloads in the JSON record):
\begin{itemize}
\item \textsc{analytical}, \texttt{analytical\_\allowbreak{}formula}; data split: n/a (analytical). Closed form from get\_\allowbreak{}formula with the on-data reading of the gate.
\item \textsc{empirical}, \texttt{empirical\_\allowbreak{}viz} (sensitivity (max\_\allowbreak{}samples 100000, sample\_\allowbreak{}fraction 1.0)); data split: train. Train-split sensitivity, activity, outcome rates and ablation.
\end{itemize}

\textbf{Authorship:} drafted by AI-assisted draft on 2026-11-10; verified by nobody (unverified).

\subsection{Audit against the definitions}\label{app:heart-s2-c3:audit}
Automatic findings on the stored explanation (codes as in the JSON record's \texttt{issues}); they flag what an examiner should check, not verdicts.
\begin{itemize}
\item \textsc{warning} \texttt{unverified}: No human verification is recorded for this explanation.
\end{itemize}
\end{drafted}
